% Artigo no modelo SBC. Compilar com o template oficial (sbc-template.sty e sbc.bst),
% disponível no site da SBC e como modelo "SBC Conferences Template" no Overleaf.
\documentclass[12pt]{article}
\usepackage{sbc-template}
\usepackage{graphicx,url}
\usepackage[utf8]{inputenc}
\usepackage[T1]{fontenc}
\usepackage[brazil]{babel}

\sloppy

\title{JuriGraph: GraphRAG e KAG sobre Neo4j\\para pesquisa em legislação brasileira}

\author{Nome Integrante 1\inst{1}, Nome Integrante 2\inst{1}}

\address{Curso de Ciência da Computação -- Instituição\\
  Cidade -- UF -- Brasil
  \email{\{integrante1,integrante2\}@instituicao.edu.br}
}

\begin{document}

\maketitle

\begin{abstract}
We present JuriGraph, a legal research assistant over Brazilian legislation that combines lexical, vector and
graph retrieval on a labeled property graph in Neo4j. Documents are converted into a knowledge graph by a
deterministic structural layer and an LLM-based semantic layer, followed by entity resolution. Answers are
produced by three retrieval strategies (vector RAG, hybrid GraphRAG with two-hop expansion and Text-to-Cypher,
and a KAG-style planner), share a citation-aware generator and a verifier, and expose an explainable confidence
score. On 32 questions written from the corpus, all strategies refused every unanswerable question, while graph
retrieval did not outperform the vector baseline in accuracy.
\end{abstract}

\begin{resumo}
Apresentamos o JuriGraph, um assistente de pesquisa sobre legislação brasileira que combina recuperação
lexical, vetorial e por travessia de grafo sobre um grafo de propriedades rotulado no Neo4j. Os documentos
viram um grafo de conhecimento por uma camada estrutural determinística e uma camada semântica com LLM,
seguidas de resolução de entidades. Três estratégias de recuperação (RAG vetorial, GraphRAG híbrido com
expansão de dois saltos e Text-to-Cypher, e um planejador no estilo KAG) compartilham um gerador com
citações e um verificador, e expõem um escore de confiança explicável. Em 32 perguntas escritas a partir do corpus,
todas as estratégias recusaram as perguntas sem resposta, e o GraphRAG híbrido teve o maior acerto
(83,3\% contra 74,1\% do baseline vetorial), ao custo de cerca de três vezes mais por pergunta.
\end{resumo}

\section{Introdução}

A pesquisa jurídica depende de relações explícitas entre normas: um decreto regulamenta uma lei, uma lei altera
artigos de outra, um tema de repercussão geral interpreta um dispositivo. O RAG tradicional
\cite{lewis2020rag} recupera trechos por similaridade e perde essas relações quando a resposta exige combinar
mais de um documento. O GraphRAG \cite{edge2024graphrag} e o KAG \cite{liang2024kag} propõem representar o
corpus como grafo e raciocinar sobre ele.

Este trabalho constrói e avalia um sistema completo com quatro pilares: (i) conversão de texto em grafo;
(ii) modelagem LPG no Neo4j com restrições e índices; (iii) recuperação híbrida com multi-hop e algoritmos de
grafo; (iv) interface que mostra o subgrafo usado e a origem de cada afirmação.

\section{Trabalhos relacionados}

O GraphRAG da Microsoft \cite{edge2024graphrag} detecta comunidades com Leiden \cite{traag2019leiden} e gera
resumos que respondem perguntas globais sobre o corpus. O KAG \cite{liang2024kag} decompõe a pergunta em um
plano lógico executado sobre o grafo e mistura passos de recuperação e de raciocínio. O HybridRAG
\cite{sarmah2024hybridrag} combina recuperação vetorial e por grafo e mostra ganhos sobre cada uma isolada.
Bancos de grafo de propriedades \cite{robinson2015graphdb} permitem guardar a evidência de cada relação, o que
usamos para rastreabilidade e para mitigar alucinação \cite{ji2023hallucination}.

\section{Arquitetura}

\subsection{Corpus e modelo de grafo}

O corpus tem 50 documentos públicos: a LGPD, o Marco Civil da Internet, o CDC, a LAI, a Lei 14.181/2021,
os Decretos 8.771/2016 e 7.962/2013 (divididos por capítulo) e cinco temas de repercussão geral do STF. O
modelo LPG tem nós \texttt{Document} e \texttt{Chunk} (camada lexical) e dez tipos de entidade, todos com o
rótulo adicional \texttt{Entity} e uma chave canônica: \texttt{Lei}, \texttt{Artigo}, \texttt{Precedente},
\texttt{Tribunal}, \texttt{Orgao}, \texttt{Conceito}, \texttt{Sujeito}, \texttt{Direito}, \texttt{Obrigacao} e
\texttt{Sancao}. Há treze tipos de relação semântica (por exemplo \texttt{REGULAMENTA}, \texttt{ALTERA},
\texttt{INTERPRETA}, \texttt{IMPOE}), cada uma com identificador, trechos de origem, camada, confiança e
evidência textual. O esquema tem restrições de unicidade, índices \emph{range}, dois índices \emph{fulltext} com
analisador para português e três índices vetoriais (trechos, entidades e comunidades).

Após a ingestão, o grafo tem 451 trechos, 2.056 entidades (732 conceitos, 427 artigos, 298 obrigações) e 8.342
relações, contando as lexicais e as de pertença a comunidade. Delas, 1.492 são semânticas, extraídas pelo LLM, e 828 são estruturais, obtidas por regras.

\subsection{Text-to-Graph}

A ingestão tem três etapas. A camada estrutural usa expressões regulares para identificar normas, artigos,
remissões (\texttt{CITA}), alterações (\texttt{ALTERA}, \texttt{REVOGA}) e regulamentações com confiança 1,0.
Ela ignora as notas de redação, que de outro modo gerariam citações falsas. A camada semântica usa o
\texttt{LLMEntityRelationExtractor} do \texttt{neo4j-graphrag} com um esquema restrito a padrões
(origem, relação, destino) permitidos. Triplas fora do esquema são descartadas. A resolução de entidades
combina chave exata, similaridade \emph{fuzzy} e similaridade de \emph{embeddings} \cite{reimers2019sbert}. Ela
tem travas contra fusões erradas (especialização, troca de termo, números diferentes e exigência de apoio
léxico) e registra cada fusão em log; na rodada final foram 52 fusões. A ingestão é idempotente, porque um
\emph{hash} por documento permite pular o que não mudou.

\subsection{Recuperação}

O modo \textbf{vetorial} é o baseline: os $k$ trechos mais próximos no índice vetorial. O modo \textbf{híbrido}
usa o \texttt{HybridCypherRetriever}, que soma vetor e BM25 com um ranqueador linear ($\alpha = 0{,}6$). A partir de
cada trecho, uma consulta Cypher expande dois saltos: das entidades mencionadas (ordenadas por PageRank
\cite{page1999pagerank}) aos vizinhos e aos vizinhos destes. Esse modo também roda um passo de Text-to-Cypher e
consulta os resumos de comunidade. O modo \textbf{KAG} pede ao LLM um plano de passos tipados (busca, Cypher,
caminho mínimo, comunidades) com dependências, e os executa em ordem topológica. O resultado de um passo pode
ser usado no seguinte.

O Text-to-Cypher recebe o esquema lido do banco e oito exemplos do domínio. Antes de executar, a consulta passa
por um validador: palavras de escrita bloqueadas fora de literais e comentários, procedimentos só de uma
lista permitida e \texttt{LIMIT} obrigatório. Depois vêm um \texttt{EXPLAIN}, a execução em sessão de leitura com
tempo-limite e uma nova tentativa com a mensagem de erro, ou quando o resultado volta vazio.

Os algoritmos de grafo rodam no GDS: PageRank e Leiden. O Leiden encontrou 532 comunidades, com modularidade
0,785, e as 27 maiores foram resumidas pelo LLM. O caminho mínimo usa \texttt{shortestPath} em Cypher.

\subsection{Geração e anti-alucinação}

O gerador recebe o contexto numerado e precisa citar, em cada afirmação, o identificador de um trecho, entidade
ou relação. Quando o contexto cobre só parte da pergunta, ele responde essa parte e declara o que o corpus não
traz, em vez de recusar tudo. O verificador descarta citações a identificadores ausentes do subgrafo e pede a um juiz LLM que diga
se cada afirmação é sustentada pelos textos citados. Só as afirmações sustentadas chegam à resposta. A
confiança é $0{,}4\,s + 0{,}3\,c + 0{,}3\,v$, em que $s$ é a similaridade média dos melhores trechos, $c$ é a
fração de citações válidas e $v$ é a fração de afirmações sustentadas. Abaixo de 0,35, ou sem nenhuma afirmação
sustentada, o sistema responde que não há evidência no grafo.

\subsection{Interface}

A interface, feita em Streamlit, mostra a resposta com notas numeradas que levam a cada fonte e a confiança
decomposta. Também mostra o subgrafo interativo (\texttt{neo4j-viz}), com o tamanho dos nós dado pelo PageRank e
os elementos citados em destaque, e abas com os trechos de origem, o Cypher executado, os passos do plano, a
verificação de cada afirmação e os nós consultados.

\section{Metodologia de avaliação}

Escrevemos 32 perguntas a partir do corpus: 10 de um salto, 12 multi-hop, 5 globais e 5 impossíveis (sem
resposta no corpus). Cada pergunta tem uma resposta esperada conferida no texto. Os três modos usam o mesmo
modelo (\texttt{gpt-5.4-mini}), o mesmo gerador e o mesmo verificador, e só a recuperação muda. Um juiz LLM
classifica cada resposta como correta (1), parcial (0,5) ou incorreta (0). Nas impossíveis, o acerto é recusar.
Também medimos:
\begin{itemize}
  \item a alucinação bruta: afirmações geradas que o verificador rejeitou;
  \item a fração de perguntas impossíveis respondidas;
  \item a fidelidade das citações;
  \item a latência;
  \item o custo estimado pelos tokens consumidos.
\end{itemize}

\section{Resultados}

% Gerado por python -m jurigraph.eval.avaliar; não editar à mão.
\begin{table}[ht]\centering\small
\caption{Resultados gerais (32 perguntas por modo).}\label{tab:geral}
\begin{tabular}{lrrr}\hline
Métrica & Vetorial & Híbrido & KAG \\ \hline
Acerto (respondíveis) & 74.1\% & 72.2\% & 53.7\% \\
Acerto (impossíveis = recusa) & 100.0\% & 100.0\% & 100.0\% \\
Alucinação bruta (afirmações não sustentadas) & 17.7\% & 27.8\% & 34.4\% \\
Impossíveis respondidas & 0.0\% & 0.0\% & 0.0\% \\
Fidelidade das citações & 100.0\% & 78.4\% & 83.8\% \\
Latência média (s) & 5.0 & 6.7 & 7.9 \\
Custo médio por pergunta (US\$) & 0.0034 & 0.0116 & 0.0070 \\
\hline\end{tabular}\end{table}
\begin{table}[ht]\centering\small
\caption{Acerto por tipo de pergunta.}\label{tab:tipos}
\begin{tabular}{lrrr}\hline
Tipo & Vetorial & Híbrido & KAG \\ \hline
Um salto & 75.0\% & 75.0\% & 50.0\% \\
Multi-hop & 83.3\% & 83.3\% & 58.3\% \\
Global & 50.0\% & 40.0\% & 50.0\% \\
Impossível & 100.0\% & 100.0\% & 100.0\% \\
\hline\end{tabular}\end{table}


% Análise da execução de 18/11/2026 (data/eval/resultados.jsonl). Se a avaliação for refeita,
% revise os números citados aqui com base em resultados.tex.
\subsection{Discussão}

O GraphRAG híbrido teve o melhor acerto nas 27 perguntas respondíveis (83,3\%), contra 74,1\% do RAG vetorial
e 74,1\% do planejador KAG. A vantagem apareceu sobretudo nas perguntas de um salto (90\% contra 75\%) e nas
globais (60\% contra 50\%). Nas multi-hop, o híbrido e o KAG empataram em 87,5\%, pouco acima do vetorial
(83,3\%). A diferença menor nas multi-hop tem uma explicação no próprio corpus: ele foi dividido em um trecho
por artigo, e muitos artigos já citam o dispositivo relacionado. O texto do Tema 987, por exemplo, transcreve o
art.~19 do Marco Civil, e a busca vetorial encontra a resposta sem atravessar o grafo.

O grafo também reduziu as recusas indevidas. O vetorial recusou 11,1\% das perguntas respondíveis, o híbrido
7,4\% e o KAG 3,7\%. Com mais contexto, o gerador encontra base para responder com mais frequência. O custo é
um contexto mais ruidoso: a alucinação bruta, isto é, as afirmações geradas que o verificador rejeitou, ficou
em 20,7\% no híbrido, contra 18,7\% no vetorial e 16,2\% no KAG. O KAG respondeu mais, mas parte dessas respostas
foi incompleta. Em duas perguntas de um salto ele descreveu o que o contexto não trazia em vez de responder, o
que o juiz classificou como incorreto.

O mecanismo de anti-alucinação funcionou nos três modos. Todas as perguntas impossíveis foram recusadas, e
nenhuma afirmação sem sustentação chegou à resposta final, porque o verificador as remove. Os erros que restaram
são recusas ou respostas incompletas, e não respostas inventadas. Uma pergunta multi-hop (a remissão da LGPD à
Lei de Arbitragem) foi recusada pelos três modos, porque a busca não trouxe o inciso que a responde.

O modo híbrido custou cerca de 3,2 vezes mais por pergunta que o vetorial (US\$ 0,0126 contra US\$ 0,0039), por
causa do contexto maior e da consulta Cypher, e foi o mais lento (8,0 s em média, contra 4,5 s). A maior parte da
latência vem das chamadas ao LLM (Cypher, geração e verificação).

\textbf{Limitações.} O conjunto tem 32 perguntas, escritas pela própria equipe a partir do corpus. O juiz é o
mesmo modelo que gera as respostas. Cada configuração rodou uma única vez, sem medir a variância entre execuções,
e respostas a uma mesma pergunta variam entre execuções. A diferença de nove pontos entre o híbrido e o vetorial
corresponde a poucas perguntas e deve ser confirmada com um conjunto maior. Por fim, o modelo de embeddings é
pequeno, e perguntas muito coloquiais às vezes não recuperam o artigo certo.


\section{Conclusão}

O JuriGraph mostra que um grafo de conhecimento jurídico, construído por regras e LLM e consultado por busca
híbrida, travessia e Text-to-Cypher validado, dá respostas com citações verificáveis e um caminho auditável da
pergunta à fonte. Na avaliação, a busca híbrida com o grafo teve o maior acerto e recusou menos perguntas
respondíveis que o baseline vetorial, e nenhuma estratégia respondeu a perguntas fora do corpus. O ganho nas
perguntas multi-hop foi pequeno, o que indica que ele depende de perguntas cuja resposta esteja de fato
espalhada por vários documentos. Como trabalhos futuros, ficam: ampliar o corpus (jurisprudência do STJ, cuja coleta automática
foi bloqueada pelo portal), avaliar com especialistas em vez de um juiz LLM e reduzir a latência do verificador.

\bibliographystyle{sbc}
\bibliography{referencias}

\end{document}
